\section{Related Work}
\label{sec:related}

\subsection{Sequential Recommendation}
SASRec~\cite{kang2018self} introduced self-attention for modeling sequential user behavior,
outperforming RNN-based methods.
DIF-SASRec~\cite{xie2022difsr} decoupled item ID, attribute, and content features into
separate attention streams, enabling fine-grained preference modeling.
We adopt DIF-SASRec and replace its original ID-based item representation
with dense BGE-M3 content embeddings, eliminating the need for ID lookup
tables and enabling zero-shot generalization to unseen items.

\subsection{Multimodal Retrieval}
CLIP~\cite{radford2021clip} demonstrated strong image--text alignment through contrastive
pre-training at scale.
BGE-M3~\cite{chen2024bgem3} provides multilingual, multi-granularity dense retrieval.
Prior work on multimodal recommendation combines visual and textual
signals but rarely integrates them with behavioral graph embeddings.

\subsection{Graph-Based Collaborative Filtering}
Cleora~\cite{cleora2021} learns item embeddings via Markov-chain propagation over
co-occurrence graphs, producing stable, high-quality behavioral representations
without explicit user--item matrix factorization.
Unlike neural collaborative filtering methods, Cleora embeddings are
deterministic and require no GPU for training, making them suitable for
catalog-scale deployment.

\subsection{Hybrid Systems}
Existing hybrid recommenders typically combine one behavioral signal with one
content modality — for example, pairing graph-based co-occurrence with text
embeddings — but treat retrieval as a single unified pipeline.
None natively distinguish between active query-driven search and passive
proactive recommendation, nor fuse visual, textual, and behavioral signals
within a continuously updated user profile.
Our system addresses this gap by coupling two complementary pipelines under
a shared stateful profile that bridges both interaction modes.
